\section{基本计数原理、排列与组合}


组合计数研究的是"满足某些条件的对象有多少个"。这一节先把最基础的三种计数工具立起来：\textbf{加法原理}、\textbf{乘法原理}与\textbf{双射计数}，然后用它们数出排列、组合与多重组合数，并给出二项式系数的几条基本恒等式。后面几节（组合恒等式、卡特兰数、球与桶、容斥与抽屉）都建立在这一节的记号与语言上。

\textbf{本节用到的记号：} $[n]=\{1,2,\dots,n\}$；$|A|$ 为有限集 $A$ 的元素个数；$A\sqcup B$ 表示不交并；$A\times B$ 为笛卡尔积；$n!$ 为阶乘；$n^{\underline{k}}=n(n-1)\cdots(n-k+1)$ 为\textbf{下降阶乘}；$\binom nk$ 为二项式系数；$\operatorname{inj}([k],[n])$ 表示从 $[k]$ 到 $[n]$ 的单射全体。

\subsection{加法原理与乘法原理}


\textbf{加法原理。} 若 $A_1,\dots,A_t$ 是两两不交的有限集，则

\[
|A_1\sqcup A_2\sqcup\cdots\sqcup A_t|=|A_1|+|A_2|+\cdots+|A_t|.
\]

一句话：\textbf{分类计数，类与类不重不漏，就把各类数量相加。} 例如午饭在 20 种面食与 30 种米饭套餐中选一种，共有 $20+30=50$ 种。用之前必须确认"两类没有重叠"，否则要用后面的容斥原理。

\textbf{乘法原理。} 若 $A,B$ 是有限集，则

\[
|A\times B|=|A|\cdot |B|.
\]

更一般地，若一个对象由 $t$ 步依次选定，第 $i$ 步有 $a_i$ 种选法且各步选择互不影响，则总选择数为 $a_1a_2\cdots a_t$。例如先在 10 道素菜中选一道、再在 10 道荤菜中选一道，共有 $10\times 10=100$ 种组合。

\textbf{基本应用：函数计数。} 从 $[n]$ 到 $[m]$ 的函数共有 $m^n$ 个：$[n]$ 中每个元素独立地在 $m$ 个像中选一个。这是后面"球与桶"模型中"可区分球、可区分桶"的计数 $m^n$ 的来源。

\subsection{阶乘、排列数与单射计数}


\textbf{阶乘。} $n!=1\cdot2\cdots n$（约定 $0!=1$）。它等于 $[n]$ 的全部排列数，也就是 $[n]$ 到自身的双射个数：选第一个位置有 $n$ 种，第二个 $n-1$ 种，依次相乘。

\textbf{排列数（下降阶乘）。} 从 $n$ 个不同对象中有顺序地选出 $k$ 个，方案数为

\[
P(n,k)=n(n-1)\cdots(n-k+1)=n^{\underline{k}}=\frac{n!}{(n-k)!}\qquad(0\le k\le n).
\]

\textbf{组合意义。} 它恰好是\textbf{单射} $f:[k]\to[n]$ 的个数：第一个位置（输入 $1$）选一个像有 $n$ 种，第二个位置（输入 $2$）不能与前面重复，剩 $n-1$ 种，依此类推。所以

\[
\#\operatorname{inj}([k],[n])=n^{\underline{k}}=\frac{n!}{(n-k)!}.
\]

例如 $P(5,3)=5\cdot4\cdot3=60$。当 $k>n$ 时约定 $n^{\underline{k}}=0$（放不下），于是公式对所有 $k\ge0$ 都成立。

\subsection{组合数与二项式系数}


\textbf{定义。} 从 $n$ 个不同对象中\textbf{不重复、不计顺序}地选出 $k$ 个，方案数记作

\[
\binom nk=\frac{n^{\underline{k}}}{k!}=\frac{n!}{k!\,(n-k)!}\qquad(0\le k\le n),
\]

称为\textbf{二项式系数}（也记 $C(n,k)$、$C_n^k$）。它计数的是 $[n]$ 的 $k$ 元子集：先按顺序选（$n^{\underline{k}}$ 种），再除掉 $k$ 个元素的 $k!$ 种排列。于是

\[
P(n,k)=\binom nk\,k! .
\]

约定：$k<0$ 或 $k>n$ 时 $\binom nk=0$；$\binom n0=\binom nn=1$。

\textbf{组合数与杨辉三角。} 把 $\binom nk$ 按 $n$ 排成三角形，就是帕斯卡三角：每一行两端是 $1$，中间每个数等于上一行左右两个数之和——这正是下面的帕斯卡恒等式。

\subsection{二项式系数的四条基本恒等式}


\textbf{恒等式一（对称性）。}

\[
\binom nk=\binom n{n-k}.
\]

\textbf{证明（双射）。} $[n]$ 的 $k$ 元子集 $S$ 与其余集 $[n]\setminus S$（$n-k$ 元子集）一一对应。$\square$

\textbf{恒等式二（帕斯卡恒等式）。} 对 $1\le k\le n-1$，

\[
\binom nk=\binom{n-1}{k-1}+\binom{n-1}{k}.
\]

\textbf{证明（分类）。} 看 $[n]$ 的 $k$ 元子集是否包含元素 $n$：包含 $n$ 的，相当于在 $[n-1]$ 中再选 $k-1$ 个，共 $\binom{n-1}{k-1}$ 种；不包含 $n$ 的，相当于在 $[n-1]$ 中选 $k$ 个，共 $\binom{n-1}{k}$ 种。两类不交且覆盖全部，由加法原理得证。$\square$

\textbf{恒等式三（吸收恒等式）。} 对 $1\le k\le n$，

\[
k\binom nk=n\binom{n-1}{k-1}.
\]

\textbf{证明（双计数）。} 左边数的是"先选 $k$ 元子集、再从其中指定一个元素"的方案数：先选子集 $\binom nk$ 种，再在 $k$ 个元素中选一个 $k$ 种。右边数的是"先指定一个元素、再从剩下的 $n-1$ 个里选 $k-1$ 个与它组成子集"：指定元素 $n$ 种，再选其余 $k-1$ 个 $\binom{n-1}{k-1}$ 种。两边数的是同一批对象（一个带标记的 $k$ 元子集）。$\square$

代数写法：$\binom nk=\frac nk\binom{n-1}{k-1}$。这条恒等式是把"$k$ 的幂"化成二项式系数的钥匙，后面求 $\sum k\binom nk$ 时反复使用。

\textbf{恒等式四（乘积恒等式）。} 对 $0\le k\le m\le n$，

\[
\binom nm\binom mk=\binom nk\binom{n-k}{m-k}.
\]

\textbf{证明（双计数）。} 左边：先从 $[n]$ 选一个 $m$ 元子集，再从中选一个 $k$ 元子集，等价于把 $[n]$ 分成三组（大小 $k$、$m-k$、$n-m$），方案数 $\frac{n!}{k!(m-k)!(n-m)!}$。右边：先选 $k$ 元子集，再从剩下的 $n-k$ 个中选 $m-k$ 个，同样得到同一批"三组划分"，方案数相同。$\square$

\textbf{组合证明的两种写法。} 上面反复用了两种方法：\textbf{(1) 双射}：把待数对象与已知对象一一对应；\textbf{(2) 双计数}：把同一个集合用两种方式数出来，两个结果必须相等。做组合题时先问"这两边分别数的是什么集合"，往往比代数变形更快。

\subsection{多重组合数与多项式展开的系数}


\textbf{定义（多重组合数）。} 设 $k_1,\dots,k_t\ge0$、$k_1+\cdots+k_t=n$。把 $n$ 个不同对象分成\textbf{有序的} $t$ 组、第 $i$ 组恰有 $k_i$ 个的方案数为

\[
\binom{n}{k_1,\dots,k_t}=\frac{n!}{k_1!\,k_2!\cdots k_t!}.
\]

\textbf{证明。} 依次从剩余对象中选：先选第 1 组 $\binom n{k_1}$ 种，再从剩下的 $n-k_1$ 个中选第 2 组 $\binom{n-k_1}{k_2}$ 种，依此类推，乘积化简即得。特别地，$t=2$ 时 $\binom{n}{k,n-k}=\binom nk$。$\square$

\textbf{二项式定理。} 对任意 $x,y$ 与 $n\ge0$，

\[
(x+y)^n=\sum_{k=0}^{n}\binom nk x^k y^{n-k}.
\]

\textbf{证明（组合意义）。} 把 $(x+y)^n$ 看成 $n$ 个因子 $(x+y)$ 相乘。展开时每个因子贡献一个 $x$ 或一个 $y$；若恰有 $k$ 个因子贡献 $x$，则得到单项式 $x^ky^{n-k}$。选择哪 $k$ 个因子贡献 $x$ 有 $\binom nk$ 种，故该项系数为 $\binom nk$。$\square$

（也可以用帕斯卡恒等式对 $n$ 归纳证明：从 $n-1$ 到 $n$ 时，$x^ky^{n-k}$ 的系数由上一行的 $\binom{n-1}{k-1}$ 与 $\binom{n-1}{k}$ 相加得到。）

\textbf{多项式定理。} 类似地，

\[
(x_1+x_2+\cdots+x_t)^n=\sum_{\substack{k_1,\dots,k_t\ge0\\ k_1+\cdots+k_t=n}}\binom{n}{k_1,\dots,k_t}x_1^{k_1}x_2^{k_2}\cdots x_t^{k_t}.
\]

例如 $(x+y+z)^n=\sum_{a+b+c=n}\frac{n!}{a!b!c!}x^ay^bz^c$。

\textbf{两个常用特值。} 在二项式定理中取 $x=y=1$ 得

\[
\sum_{k=0}^{n}\binom nk=2^n,
\]

即 $[n]$ 的全部子集个数；取 $x=1,y=-1$（$n\ge1$）得

\[
\sum_{k=0}^{n}(-1)^k\binom nk=0,
\]

即偶数元子集与奇数元子集一样多。

\subsection{本节的勘误与提醒}


\begin{itemize}
\item \textbf{加法原理要求"不交"。} 若两类有重叠，直接相加会重复计数，应改用容斥原理。
\item \textbf{排列数 vs 组合数。} 计顺序用 $P(n,k)$（单射），不计顺序用 $\binom nk$（子集）；两者差一个 $k!$。
\item \textbf{$\binom nk$ 的约定要统一。} 本节取 $k<0$ 或 $k>n$ 时为 $0$（$n\ge0$ 时）；第 2 节把 $n$ 推广到负数时，$k>n$ 不再自动为 $0$，要按下降阶乘的定义算。
\item \textbf{多重组合数是有序分组。} $\binom{n}{k_1,\dots,k_t}$ 中组的编号是固定的；若组不标号，要再除以组的对称性，那属于另一类问题（分拆数）。
\item 课堂转写把"下降阶乘/排列数"的记号念成"$n$ 的 $k$ 次方"或"$A_n^k$"，以本节的 $n^{\underline{k}}=\frac{n!}{(n-k)!}$ 为准。
\end{itemize}
